% Sources: README "Model Physics" (intro), "Units and radii", "Wind scenarios and the accretion-rate profile"
% Code: ULXspectra, branch fix-results-integrity, commit 5f00569

\section{Model overview}
\label{sec:model}

\subsection{Geometry and order of the calculation}
\label{sec:model:geometry}

The accretor is a black hole (BH) or a neutron star (NS) with a magnetosphere.
Matter arrives at the outer disc radius at a constant rate $\Mdot_0$.
Inside the spherization radius $\Rsph$ part of it leaves the disc in a wind, so that the accretion rate $\Mdot(R)$ falls towards the inner disc radius $\rin$.
The wind fills the cone between the funnel, the innermost streamline at angle $\theta_{\rm c}$ to the axis, and the disc surface.
For a NS, the matter inside $\rin = \rmag$ radiates from the magnetosphere, which we call the central emitter.
% OUTLINE: sketch figure of the geometry (disc, wind cone, funnel, photosphere,
%          outer wall, central emitter, observer at inclination i). AUTHOR: make or
%          provide; it must come from a script in the code repository.

The calculation goes through these steps in order:
\begin{enumerate}
\item radii and units (this section);
\item $\Mdot(R)$, $\Rsph$ and the trapping radius $\Rtrap$ (Sects.~\ref{sec:model:mdot} and~\ref{sec:radii});
\item viscous torque and disc effective temperature (Sect.~\ref{sec:disc});
\item luminosity budget of the disc and the central emitter (Sect.~\ref{sec:disc:budget});
\item the NS magnetosphere (Sect.~\ref{sec:central});
\item wind streamlines, density and photosphere (Sect.~\ref{sec:wind});
\item wind temperature and its energy budget (Sect.~\ref{sec:wtemp});
\item visibility and spectrum (Sect.~\ref{sec:spectrum}).
\end{enumerate}
% Source: README "Model Physics", steps 1-8.

\subsection{Units and radii}
\label{sec:model:units}

The Eddington luminosity is
\begin{equation}
\LEdd = \frac{4\pi G M c}{\kT}, \qquad \kT = 0.34~{\rm cm^2\,g^{-1}},
\label{eq:LEdd}
\end{equation}
% code: README "Units and radii"; KAPPA_T_CGS in csrcs
and the Eddington accretion rate is
\begin{equation}
\MEdd = \frac{\LEdd}{\etaEdd\, c^2}, \qquad \mdot = \frac{\Mdot}{\MEdd}.
\label{eq:MEdd}
\end{equation}
% code: README "Units and radii"; eta_edd_of() (csrcs/disc_struct.c)
We take $\etaEdd = 1/12$, that is $\MEdd = 12\,\LEdd/c^2$, as in \citet{Poutanen2007} and \citet{Lipunova1999}; \citet{Middleton2026} use $\etaEdd = 1/10$.
The outer accretion rate is $\mdot_0$.
% AUTHOR: one sentence on why 1/12 (e.g. efficiency of a Newtonian disc with
%         r_in = 6 r_g)? The sources only say "as in P07 and L99".

Radii are measured in $\rg = GM/c^2$, and in the code in units of the inner disc radius, $x = R/\rin$.
For a BH, $\rin = r_{\rm in,f}\,\rg$, where $r_{\rm in,f}$ (\code{r\_in\_factor}) is a free parameter.
For a NS, the inner disc radius is the magnetospheric radius,
\begin{equation}
\rin = \rmag = x_{\rm m}\, r_{\rm in,f}\, \rg ,
\label{eq:rm}
\end{equation}
% code: apply_physical_radii() (csrcs/disc_struct.c); README "Units and radii"
where $x_{\rm m}$ is \code{wind\_inner\_launch\_x}; only the product of the two factors matters.
The prior rejects $\rmag \le \RNS = 10$~km.
In both cases the wind is launched from $\rin$ outwards.
% AUTHOR: the code names (wind_inner_launch_x, r_in_factor) are historical; is the
%         paper notation r_in,f and x_m fine, or would you parametrise the NS by r_m
%         directly?

The flow rules of \citet{Poutanen2007} and \citet{Shakura1973} (Sect.~\ref{sec:model:mdot}) are fixed in units of $\rg$ through the reference radius $\Rref = 6\,\rg$, the innermost stable circular orbit of a non-rotating BH and the inner radius $3R_{\rm S}$ of \citet{Poutanen2007}.
The actual inner disc radius only truncates them.
% code: R_IN_REF_FACTOR (csrcs); README "Units and radii"

The outer disc radius $\rout$ is free.
Only the disc inside $\rout$ emits and counts in the disc luminosity.
If $\rout$ lies inside the outer wall of the wind (Sect.~\ref{sec:wind:photosphere}), the disc there is hidden but still sets the wind temperature (Sect.~\ref{sec:wtemp}); beyond the wall the outer disc is seen directly.
The radial grid extends to $\max(10^4\mdot_0\,\rin, \rout)$; the rings beyond $\rout$ do not emit but can occult.
% code: README "Units and radii"; radial_grid_edges() (csrcs/wind_surface.c)

\subsection{Accretion rate in the disc}
\label{sec:model:mdot}

Outside $\Rsph$ the accretion rate is constant, $\mdot = \mdot_0$.
For $\mdot_0 \le 1$ there is no wind and $\Rsph = \rin$ (Sect.~\ref{sec:radii}).
Inside $\Rsph$ the main model of this paper prescribes the fraction $f_{\rm in} = \mdot(\rin)/\mdot_0$ (\code{user\_mdot\_in\_fraction}) that reaches the inner disc radius, and takes $\Mdot$ linear in $R$ between $\rin$ and $\Rsph$:
\begin{equation}
\mdot(x) = \mdot_0 \left[ f_{\rm in} + (1 - f_{\rm in})\,\frac{x - 1}{\xsp - 1} \right]
\label{eq:mdot_setin}
\end{equation}
% code: mdot_acc_profile_scalar(), branch "set_mdot_in" (csrcs/disc_struct.c); README "Wind scenarios"
for $1 \le x < \xsp$, with $\xsp = \Rsph/\rin$.
$f_{\rm in}$ is free in the fits.
% AUTHOR: motivation for the linear profile and for leaving f_in free (vs P07 eq. 23,
%         which fixes it through eps_w)?

We compare this model with three flow rules from the literature, selected in the code by \code{wind\_scenario} (Table~\ref{tab:presets}).
In the profile of \citet[][hereafter P07]{Poutanen2007}, $\Mdot$ is linear in $u = R/\Rref$,
\begin{equation}
\mdot = \mdot_{\rm in} + (\mdot_0 - \mdot_{\rm in})\,\frac{u - 1}{u_{\rm sph} - 1}, \qquad u < u_{\rm sph},
\label{eq:mdot_P07}
\end{equation}
% code: mdot_acc_profile_scalar(), branch "P07" (csrcs/disc_struct.c); README "Wind scenarios"
with the inner rate set by the wind energy fraction $\epsilon_{\rm w}$ (P07, eq.~23),
\begin{equation}
\frac{\mdot_{\rm in}}{\mdot_0} = \frac{1 - a}{1 - a/\sqrt{2\mdot_0/5}},
\label{eq:mdot_in_P07}
\end{equation}
% code: mdot_in_over_mdot0(), a_ew() (csrcs/disc_struct.c); README "Wind scenarios"
where $a = \epsilon_{\rm w}(0.83 - 0.25\,\epsilon_{\rm w})$.
The profile is evaluated at $\max(R, \rin)$ and is not allowed to become negative.
The L99 preset is the P07 profile with $\epsilon_{\rm w} = 1$ \citep{Lipunova1999}.
In the profile of \citet[][SS73]{Shakura1973},
\begin{equation}
\mdot = \mdot_0\,\frac{R}{\Rsph}, \qquad R < \Rsph .
\label{eq:mdot_SS73}
\end{equation}
% code: mdot_acc_profile_scalar(), branch "SS73"; mdot_in_over_mdot0() gives x_launch/x_sp
% AUTHOR: check that the P07 eq. 23 and the attribution of eq. (mdot_SS73) to SS73
%         are as you want them cited.

\begin{table*}
\centering
\caption{Flow rules (\code{wind\_scenario}). The presets of the literature fix the $\Mdot(R)$ profile, the $\Rsph$ rule and $\etaEdd$; $\Rtrap$ can be set in any of them.}
\label{tab:presets}
\small
\begin{tabular}{@{}lllll@{}}
\toprule
Preset & $\Mdot(R)$ & $\Rsph$ & $\etaEdd$ & $\Rtrap$ \\
\midrule
main (\code{set\_mdot\_in}) & Eq.~\eqref{eq:mdot_setin} & free & 1/12 & M26 \\
P07 & Eq.~\eqref{eq:mdot_P07} & P07fit & 1/12 & M26 \\
L99 & Eq.~\eqref{eq:mdot_P07}, $\epsilon_{\rm w}=1$ & P07fit & 1/12 & M26 \\
SS73 & Eq.~\eqref{eq:mdot_SS73} & SS73, SS73\_eq71 & 1/12 & $\Rsph$ \\
\bottomrule
\end{tabular}
% code: psrcs/flow_rules.py PAPER_PRESETS, resolve_flow_rules(); README table "Wind scenarios"
% Rules named in the R_sph and R_trap columns: Sect. radii.
\end{table*}

\subsection{Mass loss}
\label{sec:model:massloss}

The wind carries away
\begin{equation}
\Mdot_{\rm w}(R) = \Mdot(R) - \Mdot(\rin),
\label{eq:Mdotw}
\end{equation}
% code: mdot_wind_cumulative_scalar() (csrcs/wind_struct.c); README "Wind scenarios"
the mass launched between $\rin$ and $R$.
The rate of mass loss per unit area of each side of the disc at the launch radius $R_0$ is
\begin{equation}
\dot\Sigma_{\rm w}(R_0) = \frac{1}{4\pi R_0}\,\frac{{\rm d}\Mdot_{\rm w}}{{\rm d}R_0} .
\label{eq:Sigmadotw}
\end{equation}
% code: make_dotSigma_table() (csrcs/wind_struct.c); README "Wind scenarios"
For the main model, $\dot\Sigma_{\rm w} \propto 1/R_0$ between $\rin$ and $\Rsph$.

\subsection{Parameters}
\label{sec:model:params}

Table~\ref{tab:params} lists the parameters of the model.
% OUTLINE: say which parameters are free in which fit (BH1/BH3/BH4/NS1 presets of
%          ULX.conf), once the set of fits for the paper is fixed.

\begin{table*}
\centering
\caption{Physical parameters. Default values and prior ranges are those of \code{ULX.conf}; free parameters have flat priors in the range given.}
\label{tab:params}
\small
\begin{tabular}{@{}lllll@{}}
\toprule
Symbol & Code name & Meaning & Default & Prior \\
\midrule
$\log M/M_\odot$ & \code{log\_m\_bh} & mass of the accretor & 1.0 & (0, 4) \\
$\log \mdot_0$ & \code{log\_mdot0} & outer accretion rate, units of $\MEdd$ & 2.0 & (0, 5) \\
$f_{\rm in}$ & \code{user\_mdot\_in\_fraction} & $\mdot(\rin)/\mdot_0$, main model & 0.5 & (0, 1) \\
$r_{\rm in,f}$ & \code{r\_in\_factor} & $\rin/\rg$ (BH), $\rmag/(x_{\rm m}\rg)$ (NS) & 6 & (1, 20) \\
$\log x_{\rm m}$ & \code{log\_wind\_inner\_launch\_x} & NS: $\rmag$ in units of $r_{\rm in,f}\rg$; BH: 0 & 0 & (0, 4) \\
$\log \rout/\rg$ & \code{log\_r\_out\_factor} & outer disc radius & 5.0 & (2, 10) \\
$i$ & \code{incl\_deg} & inclination, degrees & 20 & (0, 88) \\
$\epsilon_{\rm w}$ & \code{eps\_w} & fraction of the energy spent on launching the wind & 0.5 & ($10^{-4}$, 0.999) \\
$\beta_{\rm w}$ & \code{beta\_w} & wind speed at infinity, units of $\vK(R_0)$ & 1.0 & (0.05, 3) \\
$\omf$ & \code{omega\_frac} & $\cos\theta_{\rm c} - \cos\theta_{\rm d}$, solid-angle fraction of the wind & 0.7 & (0.1, 0.85) \\
$\hd$ & \code{h\_disc} & disc surface $z/R$ & 0 & (0, 0.3) \\
$\log\eta_{\rm magn}$ & \code{log\_magn\_eta} & NS: efficiency of the central emitter & 0 & ($-5$, 0) \\
$\beta_{\rm t}$ & \code{beta\_torque\_in} & NS: torque at $\rmag$ & 1 & fixed \\
$z/r$ & \code{z\_over\_r\_trap} & disc thickness in $\Rtrap$ & 0.5 & fixed \\
\bottomrule
\end{tabular}
% code: ULX.conf [DEFAULT] at commit 5f00569; README "Physical Parameters".
% Discrete switches (not in the table): wind_scenario, rsp_rule, rtrap_rule,
% disc_luminosity_mode (1+log for BH, NS), dip, wind_velocity_law (ballistic).
\end{table*}
